% 前言与导读（第 0 部分，不编号进「讲」）
\section*{写在前面：这份讲义怎么用}
\addcontentsline{toc}{section}{写在前面：这份讲义怎么用}
\markboth{写在前面}{写在前面}

\subsection*{一句话}

Kapustin--Witten 2006 \cite{kw06} 做的是这样一件事：把四维 $\mathcal N=4$ 超对称 Yang--Mills 理论做一次特定的拓扑扭化，
再把它紧化到一条 Riemann 面上；于是一个纯四维的物理问题，在低能下变成一个以 Hitchin 模空间 $\MH$ 为目标的二维 sigma 模型。
在这个二维模型里，四维的电磁对偶（S-对偶）恰好变成二维的镜像对称，而四维的 Wilson/'t~Hooft 线算子恰好变成作用在膜上的算子。
把这些对应一条一条接起来，就得到几何 Langlands 纲领的一份「物理推导」。

这份讲义不替代原文，也不替代我们那本 219 页的解经本。它的目标只有一个：
\textbf{让你能自己翻开 \cite{kw06}，从第一页读到 §12，知道每一节在做什么。}

\subsection*{三条读法}

\begin{center}
\begin{tabular}{@{}llp{8.6cm}@{}}
\toprule
路线 & 时间 & 怎么读 \\
\midrule
快览 & 约 3 小时 & 读第 0 部分 + 第 1、4、6 讲 + 每讲的「本讲核心」与「跟读提示」，跳过全部推导。 \\
标准 & 约 3 天 & 十讲顺序读，每讲后做附录 B 的 3 道习题，再按「跟读提示」回原文对应页段。 \\
深入 & 约 2 周 & 十讲 + 每讲的自检问题全部写答案 + 原文 §2--§12 逐节读 + 第 8 讲讲到的后续文献。 \\
\bottomrule
\end{tabular}
\end{center}

\subsection*{需要什么前置}

\begin{itemize}
\item \textbf{物理}：路径积分与规范场（会写 $\int \mathcal D A\, e^{\ii S}$ 就行）、超对称的基本语言（超荷、多重态）、
      知道 $\mathcal N=4$ 超对称 Yang--Mills 是什么。拓扑扭化的技术细节第 3 讲现讲。
\item \textbf{几何}：复流形与近复结构、向量丛与联络、层与凝聚层的入门概念。
      Hitchin 模空间、超凯勒商这些第 1 讲现讲；D-模在第 7 讲现讲。
\item \textbf{不需要}：不需要先学几何 Langlands 的完整理论，也不需要懂导出代数几何。
\end{itemize}

\subsection*{记号约定}

\begin{center}
\begin{tabular}{@{}ll@{}}
\toprule
记号 & 含义 \\
\midrule
$G$ & 紧单连通单李群（主要例子 $SU(N)$）；$\gC$ 是其复化李代数 \\
$\Ldual$ & Langlands 对偶群（$\LdualC$ 是其复化）；第 2 讲解释它为什么出现在这里 \\
$C$ & 亏格 $g\ge 2$ 的紧 Riemann 面（几何 Langlands 的「曲线」） \\
$\Sigma$ & 紧化用的另一条曲面；四维理论放在 $\Sigma\times C$ 上 \\
$\MH$ & $C$ 上 $G$-Higgs 丛的模空间（Hitchin 模空间），三个复结构 $I,J,K$ \\
$\tau$ & 复耦合 $\tau=\theta/2\pi+4\pi\ii/g_{\rm YM}^2$；$\Psi$ 是扭化后的典范参数 \\
$^L\!R$、$R$ & $\Ldual$ 与 $G$ 的表示；$\mathcal O(x)$、$K_C$ 是 $C$ 上的线丛与典范丛 \\
\src{B p.148} & 原文（\cite{kw06}）的印刷页锚点；本讲义所有断言都尽量给到页 \\
\bottomrule
\end{tabular}
\end{center}

\subsection*{十讲地图}

\begin{center}
\small
\begin{tabular}{@{}clp{5.4cm}l@{}}
\toprule
讲 & 主题 & 一句话 & 原文页段 \\
\midrule
1 & 舞台 & 一条曲线、一个群、一个模空间 $\MH$ & §4, pp.38--56 \\
2 & 四维主角 & $\mathcal N=4$ SYM、S-对偶、$\Ldual$、Table 1 & §1--§2, pp.2--16 \\
3 & 扭化 & 为什么要把超对称扭掉；一族 TFT 与 $\Psi$ & §3, pp.16--38 \\
4 & 紧化 & 四维的桥塌成二维的镜像对称 & §5, pp.56--78 \\
5 & 膜与线算子 & Wilson/'t~Hooft 变成作用在膜上的算子 & §6--§8, pp.78--123 \\
6 & 核心等式 & 't~Hooft 算子 $=$ Hecke 算子 & §9--§10, pp.123--184 \\
7 & 翻译器 & 从 A-膜到 D-模，得到 Hecke 特征层 & §11, pp.184--207 \\
8 & 收获与边界 & 证明了什么、没证明什么、之后读什么 & §12, pp.207--222 \\
\bottomrule
\end{tabular}
\end{center}

\subsection*{三条诚实提醒}

\begin{enumerate}
\item \textbf{这不是数学证明。} 原文是物理论证：给出对应、给证据、给一致性检查，把「为什么应该如此」讲清楚，
      但没有证明几何 Langlands。严格的证明要等到 2024 年前后 Gaitsgory 等人的纯代数几何工作（第 8 讲会说）。
\item \textbf{只做非分歧情形。} 原文只处理 $C$ 上无分歧（unramified）的情形；分歧与野分歧是 Gukov--Witten \cite{gw06} 的主题。
\item \textbf{物理侧的很多步骤是「应该」。} 原文里 ``we expect''、``should be'' 这类措辞不是客套，
      是作者在标明「这一步只有物理直觉，没有逐条验证」。本讲义会用「原文只写应该」这样的话提醒你。
\end{enumerate}

\subsection*{和另外两份材料的关系}

我们还有两份配套材料，读者可以按需取用：
\begin{itemize}
\item 一份 219 页的\textbf{解经本}（对原文逐节深读，含四层：骨架/结构/骨血/细度，附 nLab 注释与 60 条原文问答）；
\item 一次\textbf{讨论班}的八轮报告与 112 条语料（八个人每节一人报告、三席质询、独立复算）。
这份讲义是把那两份材料里「最该先知道的东西」压成十讲——所以它比解经本浅，但更适合第一次读原文。
\end{itemize}
